\section{Objections}\label{sec:objections}

The objections are to the premises, the classification, and the scope of the conclusions.

\medskip\noindent\textbf{O1. ``Religious beliefs are not hypotheses.''}\phantomsection\label{item:O1}
On this view it is a conceptual, grammatical truth that the sense of a religious belief lies in
the role it plays in a form of life: in worship, confession, thanksgiving and the way a believer
meets the world. It does not lie in any
explanatory claim about what is. To treat belief in God as a hypothesis that might be confirmed
or disconfirmed is to misread its grammar \citep{phillips1976}.

On this view such beliefs make no claim of the kind an account makes, so P-M(a) does not cover
them and the paper has no quarrel with them. If the thesis is offered as true of what religious
believers actually mean, it is an account that goes beyond the evidence about religious practice,
and by P-M(b) it is not established.

\medskip\noindent\textbf{O2. Basic belief.}\phantomsection\label{item:O2}
On Reformed epistemology, belief in God can be properly basic: formed directly by a cognitive
faculty, not inferred from other beliefs. It is warranted if that faculty functions properly, in
an appropriate environment, according to a design plan aimed at truth. If theism is true, such
belief is probably warranted, and its warrant owes nothing to argument \citep{plantinga2000}.

The paper is silent on rationality, proper function and warrant. By P-M(a), a faculty supplies a
judgment, and ``this faculty is reliable'' is an account, so by P-M(b) it is not established. On externalist theories of knowledge, the belief
may count as knowledge. Whether it does turns on whether the worldview is true, and C1 says that
cannot be established.

\medskip\noindent\textbf{O3. Faith as certain knowledge.}\phantomsection\label{item:O3}
Faith, on this view, is more certain than wisdom, science or understanding in respect of its
cause, because it is founded on the divine truth \citep[II-II q.4 a.8]{aquinas-st}. It is a firm and
certain knowledge of God's benevolence toward us \citep[III.2.7]{calvin1559}. Its assent is given
on the authority of God who reveals, who ``can neither deceive nor be deceived''
\citep[ch.~3]{vatican1870}.

By P-M(a), faith offered as firmness of assent is an imposed constraint: a choice. Faith offered as
knowledge rests on grounds that are accounts, and by P-M(b) they are not established.

\medskip\noindent\textbf{O4. Factive access and unfakeable marks.}\phantomsection\label{item:O4}
Epistemological disjunctivism holds that in the good case a person's reason for belief can be
factive and reflectively available at once: the reason is that one sees that p, not that it
seems to one that p \citep{pritchard2012}. Some traditions hold that certain experiences of God
carry a mark that cannot be counterfeited. Consolation without preceding cause, on Ignatius's
rules, is given by God alone \citep[\S330]{ignatius-exercises}.

By P-M(a) and AA-D, what such an experience supplies is evidence: how things seem. The claim that this seeming is a factive reason, or is
uncounterfeitable, is an account. By P-M(b) it has a consistent rival. Divine veracity is part of
the account, so it cannot establish the account.

\medskip\noindent\textbf{O5. ``God is beyond all categories.''}\phantomsection\label{item:O5}
On apophatic theology, God is beyond every assertion and every denial, so no positive claim about
God is true as stated \citep{dionysius-mt}.

Then the worldview asserts nothing that could be established. Once a positive claim is made, C1
applies.

\medskip\noindent\textbf{O6. ``Only God knows.''}\phantomsection\label{item:O6}
The objection is that God knows which account is true, so ``no one'' is too strong.

Grant that, if God exists, God knows. ``No one'' in the core claim means no one who reasons from
finite evidence.

\medskip\noindent\textbf{O7. ``The standard is too high.''}\phantomsection\label{item:O7}
This is the triviality objection. If establishing a claim means excluding every consistent rival,
no empirical claim is ever established. The core claim then says nothing special about religion.

Concede the parity. By P-M(b) and D3, no account that goes beyond the evidence is established,
including ordinary empirical ones. But
the standard is the one the classic project claims to meet. God, on that project, ``can be known
with certainty'' by the natural light of reason from created things \citep[ch.~2]{vatican1870},
and some critics claim a demonstration of the denial \citep{mackie1955,findlay1948}.

\medskip\noindent\textbf{O8. Atheism as absence of belief.}\phantomsection\label{item:O8}
On this view atheism is not a claim but the lack of one, and so it needs no defense.

Absence of belief makes no claim, so C1 does not apply to it. The denial that God exists is an
account, so C1 applies to it.

\medskip\noindent\textbf{O9. \emph{A priori} premises.}\phantomsection\label{item:O9}
A defender may hold that a premise classed as an imposed constraint is really a truth of reason.
An example is the possibility of a maximally great being, which is rationally acceptable on its
defender's view \citep{plantinga1974}.

If it were a truth of reason, and it excluded every rival, P-M(b) would be false for that
worldview, and C1 would not reach it.

\medskip\noindent\textbf{O10. ``The mapping begs the question.''}\phantomsection\label{item:O10}
The objection is that P-M(b) assumes what the paper sets out to show.

P-M(b) is an owned premise, with its refutation condition (\cref{item:PM}, with its motivation). The classic \emph{a priori} arguments are
attempts to meet its refutation condition, and \cref{sec:arguments} names, for each argument, the
element it would have to show to be a truth of reason (O9).
